\origpage{42}{ROBERTSON AND SEYMOUR}
\noindent result has a converse, which is a prototype of our main theorem, and which we discuss in Section 2.) Similarly, a society which can be drawn in the above way except for one crossing may have a crooked transaction of cardinality two, but has none of cardinality three or more. In contrast, a society with two crossings, sufficiently far apart, may admit a large crooked transaction. But in general, if a society can be drawn in a disc so that all its crossings are pairwise close together (in a sense which we shall specify in a moment) then it has no large crooked transaction. The main theorem of this paper is a converse; that such a society is (modulo ($\le 3$)-separations) the most general without a large crooked transaction.

Let us make the above ideas precise. We say that $(G, \Omega, \Omega_0)$ is a \textit{neighborhood} if $G$ is a graph and $\Omega, \Omega_0$ are both cyclic permutations in $V(G)$. Let $\Sigma$ be a plane, with some orientation called ``clockwise.'' We say that a neighborhood $(G, \Omega, \Omega_0)$ is \textit{rural} if $G$ has a drawing $\Gamma$ in $\Sigma$ without crossings (so $G$ is planar) and there are closed discs $\Delta_0 \subseteq \Delta \subseteq \Sigma$, such that

\hspace*{12pt}(i) the drawing $\Gamma$ uses no point of $\Sigma$ outside $\Delta$, and none in the interior of $\Delta_0$, and

\hspace*{12pt}(ii) for $v \in V(G)$, the point of $\Sigma$ representing $v$ in the drawing $\Gamma$ lies in $bd(\Delta)$ (respectively, $bd(\Delta_0)$) if and only if $v \in \bar{\Omega}$ (respectively, $v \in \bar{\Omega}_0$), and the cyclic permutation of $\bar{\Omega}$ (respectively, $\bar{\Omega}_0$) obtained from the clockwise orientation of $bd(\Delta)$ (respectively, $bd(\Delta_0)$) coincides (in the natural sense) with $\Omega$ (respectively, $\Omega_0$).

A separation of a graph $G$ is a pair of subgraphs $(G_1, G_2)$, where $G_1 \cup G_2 = G$ and $E(G_1 \cap G_2) = \emptyset$. A separation of a society $(G, \Omega)$ is a separation $(G_1, G_2)$ of $G$ such that $\bar{\Omega} \subseteq V(G_2)$. Let $(G_1, G_2)$ be a separation of $(G, \Omega)$, and let $\Omega_1$ be a cyclic permutation with $\bar{\Omega}_1 = V(G_1 \cap G_2)$. Then we say that $(G, \Omega)$ is the \textit{composition} of the society $(G_1, \Omega_1)$ with the neighborhood $(G_2, \Omega, \Omega_1)$. The \textit{sophistication} of $(G, \Omega)$ is the minimum $p \ge 0$ such that $(G, \Omega)$ is the composition of a society which has no transaction of cardinality $> p$ with a rural neighborhood. It is easy to prove (directly, or from (9.1)) that if $(G, \Omega)$ has sophistication $\le p$ then it has no crooked transaction of cardinality $> p$. The converse, however, is false. Let $G$ be the graph obtained from a cylinder graph with $3p$ radial lines and a large number of nested circles by adding three crossing pairs of edges as in Fig. 1. Let $\Omega$ be the cyclic permutation of the vertices of $G$ on the outside circle. It can be shown that $(G, \Omega)$ has no crooked transaction of cardinality $> p + 2$, and yet its sophistication is at least $\lfloor \frac{3}{2} p \rfloor$.

The main result of this paper is nevertheless a partial converse. We show that modulo ($\le 3$)-separations, a society in which there is no crooked transaction of cardinality $\ge p (\ge 4)$ has sophistication $\le 3p + 9$.
